\section{ Checking if H is Monotone}




As the $\lambda(r_j^n)^2$ is nonnegative and $\sum_{k=-1}^m c_{j,k}^n = 1$, it remains to show that $c_{j,k}^n \geq 0$ for all $k= -1, 0,1, \cdots, m.$

Clearly, $c_{j,-1}^n \geq 0$ because we have shown that in \eqref{eq:First_argument}
\begin{align*}
    c_{j,-1}^n & =\frac{\partial \mathcal{H}}{\partial \rho_{j-1}^n}              \\
               & = \lambda(1-q_j^n) \geq 0 \qquad \text{as } 0 \leq q_j^n \leq 1.
\end{align*}
For $k= 0$, we have
\begin{align*}
    c_{j,0}^n & = 1 - \lambda (1- q_j^n) + \lambda p_{j,0}^n -2\lambda r_j^n                            \\
              & = 1 - \lambda (1- q_j^n) + \lambda \Bigl(-w_0\rho_j^n + w_0 r_j^n\Bigr) -2\lambda r_j^n \\
              & = 1- \lambda \Bigl[(1- q_j^n)+ w_0\rho_j^n  + (2-w_0)r_j^n\Bigr]
\end{align*}
In order to have  $c_{j,0}^n \geq 0$, it holds to show that $\lambda \Bigl[(1- q_j^n)+ w_0\rho_j^n  + (2-w_0)r_j^n\Bigr] \leq 1$. But by \eqref{eq:lambda_God} we have $\lambda < \frac{1}{2}$. Therefore, it is sufficient to show that
\[
    (1- q_j^n)+ w_0\rho_j^n  + (2-w_0)r_j^n \leq 2.
\]
Mark that all the terms $(1- q_j^n)$, $w_0\rho_j^n$ and $(2-w_0)$ are nonnegative since we have $ 0\leq q_j^n \leq 1$,  $0 \leq  w_0 \leq 1$ and  $0\leq \rho_j^n \leq 1$. We assume first that $r_j^n \leq 0$, then we $(2-w_0)r_j^n \leq 0$ and
\begin{align*}
    (1- q_j^n)+ w_0\rho_j^n  + (2-w_0)r_j^n & \leq (1- q_j^n)+ w_0\rho_j^n \\
                                            & \leq (1- q_j^n)+ \rho_j^n    \\
                                            & \leq 2
\end{align*}
We next assume next that $r_j^n > 0$, in this case we have $(2-w_0)r_j^n > 0$. So, we need to show
\[
    (1- q_j^n)+ w_0\rho_j^n  + (2-w_0)r_j^n \leq 2.
\]



\begin{lemma}\label{lem:lemma_1}
    Suppose Assumptions~\ref{assm:assm_1}--\ref{assm:assm_4} hold.
    Assume in addition that
    \[
        0<\rho_{\min}\le \rho_0(x)\le 1,
        \qquad
        -\frac{\rho_0(y)-\rho_0(x)}{y-x}\le L
        \qquad \text{for all }x\ne y,
    \]
    for some constant \(L>0\). Assume also that
    \[
        0<\epsilon\le \frac{c\rho_{\min}}{2Lw(0)}
    \]
    and that the horizon is grid-aligned:
    \[
        \epsilon=m\Delta x,\qquad m\in\mathbb N.
    \]
    Let \(\{\rho_j^n\}\) be the numerical solution computed by \eqref{eq:mono_H}.
    Then
    \[
        r_j^n:=\rho_{j+1}^n-\rho_j^n\ge -L\Delta x
        \qquad \text{for all }j\in\mathbb Z,\ n\in\mathbb N.
    \]
\end{lemma}


\section{Convergence Properties of Numerical Schemes for the Local Model}
Before studying the convergence of numerical methods for the nonlocal model to local model, it is significant to study the convergence properties of numerical methods for the local model. Here is a thoughtful question to ask: if we have convergence criteria (or properties) for the numerical methods of the local model, do the numerical methods of the nonlocal model satisfy those criteria?
\subsection{Conservative Schemes}
The first property we require of the numerical schemes for the local model is conservation: no quantity should be created or destroyed over time. The entropy solution of the local model is conservative, in the sense that the integral of the solution is preserved over time. Therefore, we require the corresponding discrete quantity to be preserved as well.
\begin{definition}[Conservative scheme]
    The three-point update function $H$ given in \eqref{eq:three_point_H} is conservative if satisfies
    \[ \sum_{j}^{} u_{j}^{n+1} = \sum_{j}^{} u_{j}^{n} \quad \text{ for all } n .\]
\end{definition}
\subsection{Consistent Schemes}
The numerical schemes must be consistent with the local model \eqref{eq:local model}, i.e, they must approximate the desired PDE.
\begin{definition}[Consistency]
    The three-point update function $H$ with corresponding flux $F$ is consistent if
    \[ F(w,\dots, w) = f(w) \quad \text{for all } w \in \mathbb{R}.\]
\end{definition}
\subsection{Monotone Schemes}
\begin{definition}[Monotonicity]
    The three-point update function $H$ is non-decreasing in each of its arguments.
\end{definition}


\begin{assumption}
    \leavevmode
    \begin{enumerate}
        \item Denote by $\gamma_{ij}$, $1 \le i,j \le 4$, the second-order partial derivatives of $g$. That is, $\gamma_{ij}:= \frac{\partial^2 g}{\partial \rho_i \partial \rho_j }$, then
              \[
                  \gamma_{11}=\gamma_{12}=\gamma_{22}=0,
                  \qquad
                  \gamma_{33}=\gamma_{34}=\gamma_{44}=0,
              \]
              and
              \[
                  \gamma_{13},\gamma_{23},\gamma_{14},\gamma_{24} \le 0,
                  \qquad
                  \gamma_{13}+\gamma_{23}+\gamma_{14}+\gamma_{24} = -1.
              \]
    \end{enumerate}
\end{assumption}
\begin{assumption}
    Denote by $\theta_i$, $1 \le i \le 4$, the first-order partial derivatives of $g$, that is, $\theta_i:= \frac{\partial g}{\partial \rho_i}$. Then for any
    \[
        0 \le \rho_L,\rho_R,q_L,q_R \le 1,
    \]
    we have the following:
    \begin{enumerate}
        \item
              \begin{align*}
                   & \theta_1(q_L,q_R) \ge 0,
                  \qquad
                  \theta_2(q_L,q_R) \le 0,          \\
                   & \theta_3(\rho_L,\rho_R) \le 0,
                  \qquad
                  \theta_4(\rho_L,\rho_R) \le 0.
              \end{align*}

        \item
              \begin{align*}
                   & \theta_1(q_L,q_R)+\theta_3(\rho_L,\rho_R)
                  +2(\gamma_{13}+\gamma_{23}) \ge 0,                         \\
                   & \theta^{(2)}(q_L,q_R)-2(\gamma_{23}+\gamma_{24}) \le 0.
              \end{align*}

        \item
              \[
                  \theta_3(\rho_L,\rho_R)+\theta_4(\rho_L,\rho_R)
                  \le -\min\{\rho_L,\rho_R\}.
              \]
    \end{enumerate}
\end{assumption}
The Godunov numerical flux is given by
\[g(\rho_L,\rho_R,q_L,q_R)= \rho_L (1-q_R),\] and we compute the first order and second order partial derivatives of $g$ as follows:

\begin{align*}
    \theta_1(q_L,q_R)  := \frac{\partial g}{\partial \rho_L}, & \theta_2(q_L,q_R)  := \frac{\partial g}{\partial \rho_R} \\
    = 1-q_R                                                   & = 0
\end{align*}
\begin{align*}
    \theta_3(\rho_L,\rho_R)  := \frac{\partial g}{\partial q_L}, & \theta_4(\rho_L,\rho_R) := \frac{\partial g}{\partial q_R} \\
    = 0                                                          & = -\rho_L.
\end{align*}
Since $\theta_2$ and $\theta_3$ are zero, and  $\theta^{(1)}$ depends only on $q_R$ while $\theta^{(4)}$ on $\rho_L$, the only nonzero second order partial derivatives are the following
\begin{align*}
    \gamma_{14} & := \frac{\partial^2 g}{ \partial q_R \partial \rho_L}, & \gamma_{41} & := \frac{\partial^2 g}{\partial \rho_L \partial q_R} \\
                & = -1                                                   &             & = -1.
\end{align*}
Now we see the conditions in Assumptions 4.2.4--4.2.5 are satisfies for the Godunov numerical flux is straightforward. let us now see what does the CFL $\lambda$ in Assumptions 4.2.6 satisfy explicitly. Given
\[
    0 \le \rho_L,\rho_R,q_L,q_R \le 1,
\] as Assumption 4.2.5, we obtain
\begin{equation}\label{eq:lambda_God}
    \begin{split}
        1 > \lambda \sum_{i=1}^4 \|\theta^{(i)}\|_{L^{\infty}}                                                                                           \\
        \lambda \Bigl(\|\theta^{(1)}\|_{L^{\infty}}+ \|\theta^{(2)}\|_{L^{\infty}} + \|\theta^{(3)}\|_{L^{\infty}} + \|\theta^{(4)}\|_{L^{\infty}}\Bigr) \\
        \lambda \Bigl(\|\theta^{(1)}\|_{L^{\infty}}+ 0 + 0 + \|\theta^{(4)}\|_{L^{\infty}}\Bigr)                                                         \\
        \lambda \Bigl(\|\theta^{(1)}\|_{L^{\infty}}+ \|\theta^{(4)}\|_{L^{\infty}}\Bigr)                                                                 \\
        \lambda \Bigl( \sup\limits_{q_R \in [0,1]}| 1-q_R\vert  +  \sup\limits_{\rho_L \in [0,1]}|-\rho_L\vert \Bigr)                                    \\
        2\lambda
    \end{split}
\end{equation}
So, for the Godunov numerical flux $g$,


\begin{proof}[Proof of Lemma 4.2.9 for the Godunov flux]
    For the Godunov numerical flux
    \[
        g(\rho_L,\rho_R,q_L,q_R)=\rho_L(1-q_R),
    \]
    we have
    \[
        \gamma_{13}=\gamma_{23}=\gamma_{24}=0,
        \qquad
        \gamma_{14}=-1.
    \]
    Hence \eqref{eq:differences} reduces to
    \[
        r_j^{n+1}
        =
        \sum_{k=-1}^m c_{j,k}^n r_{j+k}^n
        +2\lambda(r_j^n)^2,
    \]
    where
    \[
        \sum_{k=-1}^m c_{j,k}^n=1,
    \]
    and
    \[
        c_{j,-1}^n=\lambda(1-q_j^n),
    \]
    \[
        c_{j,0}^n
        =
        1-\lambda(1-q_j^n)+\lambda p_{j,0}^n-2\lambda r_j^n,
    \]
    \[
        c_{j,k}^n=\lambda p_{j,k}^n,
        \qquad k=1,\dots,m,
    \]
    with
    \[
        p_{j,k}^n
        =
        w_{k-1}\rho_{j+1}^n-w_k\rho_j^n+w_k r_j^n,
        \qquad
        w_{-1}=w_m=0.
    \]

    The initial one-sided Lipschitz condition gives
    \[
        r_j^0\ge -L\Delta x,
        \qquad j\in\mathbb Z.
    \]
    Assume inductively that
    \[
        r_j^n\ge -L\Delta x,
        \qquad j\in\mathbb Z.
    \]
    We prove the same estimate at time level \(n+1\).

    First,
    \[
        c_{j,-1}^n=\lambda(1-q_j^n)\ge0,
    \]
    because \(0\le q_j^n\le1\).

    Next consider \(c_{j,0}^n\). Since
    \[
        p_{j,0}^n=-w_0\rho_j^n+w_0r_j^n,
    \]
    we have
    \[
        c_{j,0}^n
        =
        1-\lambda\Bigl[
            (1-q_j^n)+w_0\rho_j^n+(2-w_0)r_j^n
            \Bigr].
    \]
    Set
    \[
        B_j^n
        :=
        (1-q_j^n)+w_0\rho_j^n+(2-w_0)r_j^n.
    \]

    If \(r_j^n\le0\), then
    \[
        (2-w_0)r_j^n\le0.
    \]
    Moreover, since
    \[
        q_j^n
        =
        \sum_{k=0}^{m-1}w_k\rho_{j+k}^n
        \ge w_0\rho_j^n,
    \]
    we get
    \[
        B_j^n
        \le
        (1-q_j^n)+w_0\rho_j^n
        \le 1.
    \]
    Thus
    \[
        c_{j,0}^n\ge 1-\lambda>0.
    \]

    Now suppose \(r_j^n>0\). By the induction hypothesis,
    \[
        \rho_{j+k}^n
        \ge
        \rho_{j+1}^n - L(k-1)\Delta x,
        \qquad k\ge1.
    \]
    Therefore
    \[
        \begin{aligned}
            q_j^n
             & =
            w_0\rho_j^n+\sum_{k=1}^{m-1}w_k\rho_{j+k}^n \\
             & \ge
            w_0\rho_j^n
            +
            \sum_{k=1}^{m-1}w_k
            \Bigl(\rho_{j+1}^n-L(k-1)\Delta x\Bigr)     \\
             & =
            w_0\rho_j^n
            +(1-w_0)\rho_{j+1}^n
            -
            L\Delta x\sum_{k=1}^{m-1}(k-1)w_k.
        \end{aligned}
    \]
    Since \(m=\lceil \epsilon/\Delta x\rceil\), we have
    \[
        (m-1)\Delta x<\epsilon.
    \]
    Hence, for every \(1\le k\le m-1\),
    \[
        (k-1)\Delta x\le (m-2)\Delta x <\epsilon.
    \]
    Using \(\sum_{k=0}^{m-1}w_k=1\), we obtain
    \[
        \Delta x\sum_{k=1}^{m-1}(k-1)w_k
        \le \epsilon.
    \]
    Thus
    \[
        q_j^n
        \ge
        w_0\rho_j^n
        +(1-w_0)\rho_{j+1}^n
        -L\epsilon.
    \]
    Consequently,
    \[
        \begin{aligned}
            B_j^n
             & =
            (1-q_j^n)+w_0\rho_j^n+(2-w_0)r_j^n \\
             & \le
            1+L\epsilon
            -(1-w_0)\rho_{j+1}^n
            +(2-w_0)(\rho_{j+1}^n-\rho_j^n)    \\
             & =
            1+L\epsilon
            +\rho_{j+1}^n
            -(2-w_0)\rho_j^n                   \\
             & \le
            2+L\epsilon,
        \end{aligned}
    \]
    because \(0\le \rho_j^n,\rho_{j+1}^n\le1\). Therefore
    \[
        c_{j,0}^n
        \ge
        1-\lambda(2+L\epsilon).
    \]
    If
    \[
        \epsilon\le \frac{1-2\lambda}{\lambda L},
    \]
    then
    \[
        c_{j,0}^n\ge0.
    \]

    It remains to show that \(c_{j,k}^n\ge0\) for \(k=1,\dots,m\). Since
    \[
        c_{j,k}^n=\lambda p_{j,k}^n,
    \]
    it is enough to prove \(p_{j,k}^n\ge0\). We compute
    \[
        \begin{aligned}
            p_{j,k}^n
             & =
            w_{k-1}\rho_{j+1}^n-w_k\rho_j^n+w_k r_j^n \\
             & =
            (w_{k-1}-w_k)\rho_{j+1}^n+2w_k r_j^n.
        \end{aligned}
    \]

    If \(r_j^n\ge0\), then
    \[
        p_{j,k}^n
        \ge
        (w_{k-1}-w_k)\rho_{j+1}^n
        \ge
        cm^{-2}\rho_{\min}
        \ge0.
    \]

    If \(r_j^n<0\), then by the induction hypothesis,
    \[
        r_j^n\ge -L\Delta x.
    \]
    Also, from the quadrature assumption,
    \[
        w_k\le w_\epsilon(k\Delta x)\Delta x.
    \]
    Since
    \[
        w_\epsilon(s)=\frac1\epsilon w\!\left(\frac{s}{\epsilon}\right)
    \]
    and \(w\) is decreasing, we have
    \[
        w_k\le w(0)\frac{\Delta x}{\epsilon}.
    \]
    Thus
    \[
        \begin{aligned}
            p_{j,k}^n
             & \ge
            cm^{-2}\rho_{\min}
            -2w_kL\Delta x \\
             & \ge
            cm^{-2}\rho_{\min}
            -
            2Lw(0)\frac{(\Delta x)^2}{\epsilon}.
        \end{aligned}
    \]
    Since \(m=\lceil \epsilon/\Delta x\rceil\), for \(m\ge2\),
    \[
        \Delta x
        <
        \frac{\epsilon}{m-1}
        \le
        \frac{2\epsilon}{m}.
    \]
    Therefore
    \[
        \frac{(\Delta x)^2}{\epsilon}
        \le
        \frac{4\epsilon}{m^2}.
    \]
    Hence
    \[
        p_{j,k}^n
        \ge
        m^{-2}
        \Bigl(
        c\rho_{\min}-8Lw(0)\epsilon
        \Bigr).
    \]
    If
    \[
        \epsilon\le \frac{c\rho_{\min}}{8Lw(0)},
    \]
    then
    \[
        p_{j,k}^n\ge0.
    \]
    Therefore
    \[
        c_{j,k}^n\ge0,
        \qquad k=1,\dots,m.
    \]

    We have shown that
    \[
        c_{j,k}^n\ge0,
        \qquad k=-1,0,1,\dots,m,
    \]
    and
    \[
        \sum_{k=-1}^m c_{j,k}^n=1.
    \]
    Thus \(r_j^{n+1}\) is a convex combination of
    \[
        r_{j-1}^n,r_j^n,\dots,r_{j+m}^n
    \]
    plus the nonnegative term \(2\lambda(r_j^n)^2\). Therefore
    \[
        r_j^{n+1}
        \ge
        \inf_{\ell\in\mathbb Z} r_\ell^n
        \ge
        -L\Delta x.
    \]
    By induction,
    \[
        r_j^n\ge -L\Delta x,
        \qquad j\in\mathbb Z,\ n\in\mathbb N.
    \]
\end{proof}

Clearly, $c_{j,-1}^n \geq 0$ because we have shown that in \eqref{eq:First_argument}
\begin{align*}
    c_{j,-1}^n & =\frac{\partial \mathcal{H}}{\partial \rho_{j-1}^n}              \\
               & = \lambda(1-q_j^n) \geq 0 \qquad \text{as } 0 \leq q_j^n \leq 1.
\end{align*}
For $k= 0$, we have
\begin{align*}
    c_{j,0}^n & = 1 - \lambda (1- q_j^n) + \lambda p_{j,0}^n -2\lambda r_j^n                            \\
              & = 1 - \lambda (1- q_j^n) + \lambda \Bigl(-w_0\rho_j^n + w_0 r_j^n\Bigr) -2\lambda r_j^n \\
              & = 1- \lambda \Bigl[(1- q_j^n)+ w_0\rho_j^n  + (2-w_0)r_j^n\Bigr]
\end{align*}
In order to have  $c_{j,0}^n \geq 0$, it holds to show that $\lambda \Bigl[(1- q_j^n)+ w_0\rho_j^n  + (2-w_0)r_j^n\Bigr] \leq 1$. But by \eqref{eq:lambda_God} we have $\lambda < \frac{1}{2}$. Therefore, it is sufficient to show that
\[
    (1- q_j^n)+ w_0\rho_j^n  + (2-w_0)r_j^n \leq 2.
\]
Mark that all the terms $(1- q_j^n)$, $w_0\rho_j^n$ and $(2-w_0)$ are nonnegative since we have $ 0\leq q_j^n \leq 1$,  $0 \leq  w_0 \leq 1$ and  $0\leq \rho_j^n \leq 1$.

If $r_j^n \leq 0$, then we have $(2-w_0)r_j^n \leq 0$ and
\begin{align*}
    (1- q_j^n)+ w_0\rho_j^n  + (2-w_0)r_j^n & \leq (1- q_j^n)+ w_0\rho_j^n \\
                                            & \leq (1- q_j^n)+ \rho_j^n    \\
                                            & \leq 2
\end{align*}
If $r_j^n > 0$, then we have $(2-w_0)r_j^n > 0$. So, we need to show
\[
    (1- q_j^n)+ w_0\rho_j^n  + (2-w_0)r_j^n \leq 2 \qquad \text{\textcolor{red}{I NEED HELP TO COMPLETE THIS PART!}}
\]


\begin{equation}\label{eq:differences}
    r_j^{n+1}
    =
    \sum_{k=-1}^m c_{j,k}^n r_{j+k}^n
    -2\lambda(\gamma_{13}+\gamma_{14})(r_j^n)^2
    -2\lambda(\gamma_{23}+\gamma_{24})(r_{j+1}^n)^2,
\end{equation}
such that
\[\sum_{k=-1}^m c_{j,k}^n = 1.\]

Since we have
\[ \gamma_{13}= \gamma_{23}= \gamma_{24}=0 \qquad \text{and } \gamma_{14} = -1\]
for the Godunov numerical flux $g$, \eqref{eq:differences} simplifies to



It remains to prove the upper bound. Since
\[
    \mathcal H(\rho_{\max},\rho_{\max},\dots,\rho_{\max})
    =
    \rho_{\max},
\]
we write
\[
    \rho_{\max}
    -
    \mathcal H(\rho_{j-1}^n,\rho_j^n,\rho_{j+1}^n,\dots,\rho_{j+m}^n)
    =
    \Delta \mathcal K_1+\Delta \mathcal K_2,
\]
where
\[
    \begin{aligned}
        \Delta \mathcal K_1
         & =
        \mathcal H(\rho_{\max},\rho_{\max},\rho_{j+1}^n,\dots,\rho_{j+m}^n)
        -
        \mathcal H(\rho_{j-1}^n,\rho_j^n,\rho_{j+1}^n,\dots,\rho_{j+m}^n),
        \\
        \Delta \mathcal K_2
         & =
        \mathcal H(\rho_{\max},\rho_{\max},\rho_{\max},\dots,\rho_{\max})
        -
        \mathcal H(\rho_{\max},\rho_{\max},\rho_{j+1}^n,\dots,\rho_{j+m}^n).
    \end{aligned}
\]

By the mean value theorem, there exist intermediate values
\[
    \rho_{\min}\le \widehat\rho_{j-1}^n,\widehat\rho_j^n\le \rho_{\max}
\]
such that
\[
    \begin{aligned}
        \Delta \mathcal K_1
         & =
        \frac{\partial \mathcal H}{\partial \rho_{j-1}^n}
        (\widehat\rho_{j-1}^n,\widehat\rho_j^n,\rho_{j+1}^n,\dots,\rho_{j+m}^n)
        (\rho_{\max}-\rho_{j-1}^n)
        \\
         & \quad+
        \frac{\partial \mathcal H}{\partial \rho_j^n}
        (\widehat\rho_{j-1}^n,\widehat\rho_j^n,\rho_{j+1}^n,\dots,\rho_{j+m}^n)
        (\rho_{\max}-\rho_j^n).
    \end{aligned}
\]
As in the proof of the lower bound,
\[
    \frac{\partial \mathcal H}{\partial \rho_{j-1}^n}\ge 0,
    \qquad
    \frac{\partial \mathcal H}{\partial \rho_j^n}\ge 0.
\]
Moreover,
\[
    \rho_{\max}-\rho_{j-1}^n\ge 0,
    \qquad
    \rho_{\max}-\rho_j^n\ge 0.
\]
Therefore,
\[
    \Delta \mathcal K_1\ge 0.
\]

Next, by the mean value theorem, there exist intermediate values
\[
    \rho_{\min}\le \widehat\rho_{j+k}^n\le \rho_{\max},
    \qquad k=1,\dots,m,
\]
such that
\[
    \Delta \mathcal K_2
    =
    \sum_{k=1}^m
    \frac{\partial \mathcal H}{\partial \rho_{j+k}^n}
    (\rho_{\max},\rho_{\max},
    \widehat\rho_{j+1}^n,\dots,\widehat\rho_{j+m}^n)
    (\rho_{\max}-\rho_{j+k}^n).
\]
For \(k=m\), using \(w_m=0\),
\[
    \frac{\partial \mathcal H}{\partial \rho_{j+m}^n}
    =
    \lambda w_{m-1}\rho_{\max}\ge 0.
\]
For \(k=1,\dots,m-1\),
\[
    \begin{aligned}
        \frac{\partial \mathcal H}{\partial \rho_{j+k}^n}
        (\rho_{\max},\rho_{\max},
        \widehat\rho_{j+1}^n,\dots,\widehat\rho_{j+m}^n)
         & =
        \lambda(w_{k-1}\rho_{\max}-w_k\rho_{\max})
        \\
         & =
        \lambda\rho_{\max}(w_{k-1}-w_k)
        \ge 0.
    \end{aligned}
\]
Since also
\[
    \rho_{\max}-\rho_{j+k}^n\ge 0,
    \qquad k=1,\dots,m,
\]
we get
\[
    \Delta \mathcal K_2\ge 0.
\]
Hence
\[
    \rho_{\max}
    -
    \mathcal H(\rho_{j-1}^n,\rho_j^n,\dots,\rho_{j+m}^n)
    =
    \Delta \mathcal K_1+\Delta \mathcal K_2
    \ge 0,
\]
and therefore
\[
    \mathcal H(\rho_{j-1}^n,\rho_j^n,\dots,\rho_{j+m}^n)
    \le \rho_{\max}.
\]
Together with the lower bound, this proves
\[
    \rho_{\min}\le \rho_j^{n+1}\le \rho_{\max}.
\]

It remains to consider the case \(r_j^n>0\). Set
\[
    a:=L\Delta x .
\]
By the induction hypothesis, \(r_\ell^n\ge -a\) for every \(\ell\in \mathbb Z\). Hence, for every \(k=1,\dots,m-1\),
\[
    \rho_{j+k}^n
    =
    \rho_{j+1}^n+\sum_{\ell=1}^{k-1} r_{j+\ell}^n
    \ge
    \rho_{j+1}^n-(k-1)a .
\]
Therefore,
\[
    \begin{aligned}
        q_j^n
         & =
        w_0\rho_j^n+\sum_{k=1}^{m-1}w_k\rho_{j+k}^n \\
         & \ge
        w_0\rho_j^n+\sum_{k=1}^{m-1}w_k
        \bigl(\rho_{j+1}^n-(k-1)a\bigr)             \\
         & =
        w_0\rho_j^n+(1-w_0)\rho_{j+1}^n
        -
        a\sum_{k=1}^{m-1}(k-1)w_k .
    \end{aligned}
\]
Since \(\rho_{j+1}^n=\rho_j^n+r_j^n\), this gives
\[
    q_j^n
    \ge
    \rho_j^n+(1-w_0)r_j^n
    -
    a\sum_{k=1}^{m-1}(k-1)w_k .
\]
Consequently,
\[
    \begin{aligned}
         & (1-q_j^n)+w_0\rho_j^n+(2-w_0)r_j^n \\
         & \le
        1-\rho_j^n-(1-w_0)r_j^n
        +a\sum_{k=1}^{m-1}(k-1)w_k
        +w_0\rho_j^n+(2-w_0)r_j^n             \\
         & =
        1-(1-w_0)\rho_j^n+r_j^n
        +a\sum_{k=1}^{m-1}(k-1)w_k .
    \end{aligned}
\]
Because \(r_j^n>0\) and \(0\le \rho_{j+1}^n\le1\), we have
\[
    r_j^n=\rho_{j+1}^n-\rho_j^n\le 1-\rho_j^n .
\]
Thus,
\[
    \begin{aligned}
        (1-q_j^n)+w_0\rho_j^n+(2-w_0)r_j^n
         & \le
        2-(2-w_0)\rho_j^n
        +a\sum_{k=1}^{m-1}(k-1)w_k \\
         & \le
        2-(2-w_0)\rho_{\min}
        +a\sum_{k=1}^{m-1}(k-1)w_k .
    \end{aligned}
\]
We now estimate the last term. By Assumption~\ref{assm:assm_3},
\[
    w_k\le \omega^\epsilon(k\Delta x)\Delta x
    =
    \frac{1}{\epsilon}\omega\!\left(\frac{k\Delta x}{\epsilon}\right)\Delta x
    \le
    \omega(0)m^{-1}.
\]
Therefore,
\[
    \begin{aligned}
        a\sum_{k=1}^{m-1}(k-1)w_k
         & \le
        L\Delta x\,\omega(0)m^{-1}
        \sum_{k=1}^{m-1}(k-1) \\
         & \le
        \frac{L\epsilon\,\omega(0)}{2}.
    \end{aligned}
\]
Using the assumption
\[
    c\rho_{\min}\ge 2L\omega(0)\epsilon ,
\]
we get
\[
    \frac{L\epsilon\,\omega(0)}{2}
    \le
    \frac{c\rho_{\min}}{4}.
\]
Moreover, from \(w_{k-1}-w_k\ge cm^{-2}\), \(w_m=0\), and \(\sum_{k=0}^{m-1}w_k=1\), we have
\[
    1=\sum_{k=0}^{m-1}w_k
    \ge
    \sum_{k=0}^{m-1}(m-k)cm^{-2}
    =
    \frac{c(m+1)}{2m},
\]
so \(c\le 2m/(m+1)<2\). Hence
\[
    \frac{c\rho_{\min}}{4}\le \frac{\rho_{\min}}{2}.
\]
Since \(0\le w_0\le1\), we also have \(2-w_0\ge1\). Therefore,
\[
    \begin{aligned}
        (1-q_j^n)+w_0\rho_j^n+(2-w_0)r_j^n
         & \le
        2-(2-w_0)\rho_{\min}
        +\frac{\rho_{\min}}{2} \\
         & \le
        2-\frac{\rho_{\min}}{2}
        \le 2.
    \end{aligned}
\]
Thus,
\[
    \lambda\Bigl[(1-q_j^n)+w_0\rho_j^n+(2-w_0)r_j^n\Bigr]
    \le 2\lambda <1,
\]
and hence
\[
    c_{j,0}^n
    =
    1-\lambda\Bigl[(1-q_j^n)+w_0\rho_j^n+(2-w_0)r_j^n\Bigr]
    \ge 0 .
\]





the numerical solution $\mathcal{U}^{\epsilon, \Delta x}$


The well-posedness of the local conservation laws is stated in (\textcolor{red}{Citation is needed}) through viscous approximation. That means unique entropy solutions exist. Moreover the entropy solutions satisfy $L^p$ stability estimates and are Total Variation Diminishing (TVD). Accordingly in  (\textcolor{red}{Citation is needed}) they designed conservative, consistent and monotone numerical scheme, such methods converge to the unique entropy solutions of the conservation laws as they preserve the stability estimates. The well-posedness of the nonlocal conservation laws (with non-local flux arising in traffic flow modeling) is shown in (\textcolor{red}{Citation is needed}), the proof is based on a priori $L^{\infty}$ and (TVD) estimates of a numerical method.
To obtain convergence of numerical methods for local conservative laws, it suffices to design conservative, consistent and monotone numerical schemes. However, for nonlocal conservation laws it's not a straight forward. In addition as nonlocal conservation laws are meant to be generalization of the local conservation laws, one wants to obtain the so-called asymptotically compatibility, that is as the parameters tend to zero one obtains the solution of the local conservation laws. It is therefore things get complicated when designing numerical methods for nonlocal conservation laws.
They designed in (\textcolor{red}{Citation is needed}) a general class of numerical methods, that is a family of finite volume schemes. Under certain conditions such that the initial data can't have a negative jump, they proved the asymptotically compatibility and asymptotic preserving property. They provided in addition numerical experiments only for the Lax--Friedrchs type scheme. In this thesis we provide further numerical experiments where we take a comparison of the Lax--Friedrchs type scheme and Godunov type scheme, and conclude that the Godunov type scheme is the efficient one.

Those scheme are however first order accurate. High-resolution schemes are needed in practice in order to increase the computational efficiency. A second order numerical scheme for the nonlocal conservation laws have been devolved in  (\textcolor{red}{Citation is needed}).







Through out this thesis we focus on the specific one dimensional traffic flow model which is
modeled by a nonlocal conservation law. In other words the goal is to design an efficient numerical method and ultimately to prove that the method convergences to the correct solution.